% ==============================================================
%  CV TEMPLATE — Estilo "Diario Oficial"
%  Autor original: Giorgio Reveco (github.com/girebz)
%  Licencia: MIT — libre uso, atribución apreciada
% ==============================================================
%
%  INSTRUCCIONES DE USO:
%  1. Reemplaza los campos [ASI] con tu información
%  2. Comenta con % las secciones que no apliquen a tu perfil
%  3. Para adaptar a una postulación específica, comenta los
%     bloques irrelevantes sin eliminarlos
%  4. El perfil profesional y el tagline bajo tu nombre son los
%     dos elementos que más conviene personalizar por postulación
%
%  ESTRUCTURA DEL DOCUMENTO:
%  - Encabezado (nombre, tagline, contacto)
%  - Perfil profesional
%  - Habilidades
%  - Experiencia profesional (subsecciones por área)
%  - Educación y formación académica
%  - Reconocimientos
%  - Otras certificaciones
% ==============================================================

\documentclass[10pt,a4paper]{article}
\usepackage[utf8]{inputenc}
\usepackage[spanish]{babel}   % Cambiar a [english] si se prefiere
\usepackage[none]{hyphenat}
\usepackage{amsmath}
\usepackage{amsfonts}
\usepackage{amssymb}
\usepackage{graphicx}
\usepackage{ragged2e}
\usepackage{hyperref}
\usepackage{multicol,xcolor}
\usepackage[left=1.4cm,right=1.4cm,top=1cm,bottom=1cm]{geometry}
\setlength{\parindent}{0pt}
% --- PAQUETES OPCIONALES (descomentar si se requieren) --------
%\usepackage{CJKutf8}   % Para caracteres en chino/japonés/coreano
%\usepackage{background}

% --- CONFIGURACIÓN DE IMAGEN DE FONDO (opcional) --------------
%\backgroundsetup{
%  scale=1,
%  color=black,
%  opacity=0,
%  angle=0,
%  position=current page.south,
%  vshift=14.85cm,
%  hshift=0cm,
%  contents={\includegraphics[width=\paperwidth,height=\paperheight]{bg.jpg}}
%}

% --- MODO OSCURO (opcional) -----------------------------------
%\pagecolor{black}
%\color{white}

% ==============================================================
\begin{document}\sloppy
\pagenumbering{gobble}

% --- ENCABEZADO -----------------------------------------------
\noindent\rule{\textwidth}{0.5pt}\vspace{-0.25cm}
\noindent\rule{\textwidth}{2pt}\vspace{-0.4cm}
\begin{center}
  \Large{\textsc{Nombre Completo}}\\ \vspace{0.1cm}
  % Tagline: títulos, rol actual o descripción breve.
  % Consejo: adapta esta línea según la postulación.
  \large{Título 1, Título 2, Rol Actual}\\
  % Foto opcional:
  %\includegraphics[scale=0.48]{foto.jpg}\\
\end{center}
\vspace{-0.5cm}\noindent\rule{\textwidth}{0.5pt}\vspace{-0.25cm}
\noindent\rule{\textwidth}{2pt}\vspace{-0.3178cm}

% --- CONTACTO -------------------------------------------------
\begin{center}
\textbf{Email:} correo@dominio.com $|$
\textbf{LinkedIn:} \href{https://linkedin.com/in/usuario}{linkedin.com/in/usuario} $|$
\textbf{Fono:} +XX X XXXX XXXX\\

% --- PERFIL PROFESIONAL ---------------------------------------
\vspace{-0.3cm}\noindent\rule{\textwidth}{0.5pt}
\textbf{PERFIL PROFESIONAL}\vspace{-0.25cm}
\begin{justify}
% 3-5 líneas. Describe tu área de expertise, trayectoria principal
% y lo que aportas. Adaptar por postulación es altamente recomendado.
Descripción profesional: formación, trayectoria, enfoque y valor diferencial.
\end{justify}
\end{center}

% --- HABILIDADES ----------------------------------------------
\vspace{-0.10cm}\textbf{HABILIDADES}\hfill
% Enlace a portafolio o sitio web (comentar si no aplica)
\textbf{Portafolio: }\href{https://github.com/usuario}{github.com/usuario}
\vspace{-0.20cm}
\noindent\rule{\textwidth}{2pt}\phantom{a}
\\ \vspace{-0.9cm} \\
\begin{flushleft}
\textbf{Idiomas:} Idioma 1 (nivel), Idioma 2 (nivel).\\
\textbf{Programación y Software:} Lenguajes, frameworks, herramientas.\\
\textbf{Gestión y Análisis:} Metodologías, habilidades de gestión.\\
\textbf{Comunicación e Innovación:} Habilidades blandas y transversales.
\end{flushleft}\phantom{a}
\\ \vspace{-1cm} \\

% ==============================================================
%  EXPERIENCIA PROFESIONAL
% ==============================================================
\textbf{EXPERIENCIA PROFESIONAL}\vspace{-0.25cm}\\
\noindent\rule{\textwidth}{2pt}

% --- BLOQUE DE EXPERIENCIA ------------------------------------
% Duplica este bloque completo por cada área o tipo de experiencia.
% Comenta bloques enteros con % para adaptar la postulación.
% --------------------------------------------------------------

\vspace{-0.1cm}ÁREA DE EXPERIENCIA 1 \hfill (Año inicio - Presente)\vspace{-0.20cm}
\noindent\rule{\textwidth}{0.1pt}
\textbf{(Año - Presente) Rol} en la \textsc{Institución} \hfill Ciudad, País\\
\textbf{(Año - Año) Rol} en la \textsc{Institución} \hfill Ciudad, País\\
% Agrega o elimina líneas según corresponda

\vspace{0.2cm}
% Descripción del rol. Puede dividirse por sub-rol si aplica.
\textbf{Como Sub-rol o función,} descripción de responsabilidades,
metodologías aplicadas y resultados obtenidos.

\vspace{0.2cm}
% Lista de ítems en 3 columnas (útil para cursos, proyectos, etc.)
% Comentar este bloque si no aplica.
\begin{flushleft}
\begin{multicols}{3}\textit{
\hspace{-0.12cm}$\circ$ Ítem 1\\
$\circ$ Ítem 2\\
$\circ$ Ítem 3\\
$\circ$ Ítem 4\\
$\circ$ Ítem 5\\
$\circ$ Ítem 6}
\end{multicols}
\end{flushleft}

% --- SEGUNDO BLOQUE DE EXPERIENCIA ----------------------------
\vspace{-0.1cm}ÁREA DE EXPERIENCIA 2 \hfill (Año inicio - Presente)\vspace{-0.20cm}
\noindent\rule{\textwidth}{0.1pt}

\textbf{(Año - Presente) Rol} en \textsc{Institución} \hfill Ciudad, País.\\
Descripción de responsabilidades y logros concretos. Incluye métricas
cuando sea posible: porcentajes, personas impactadas, resultados medibles.

\vspace{0.2cm}
\textbf{(Año - Año) Rol} en \textsc{Institución} \hfill Ciudad, País.\\
Descripción breve del rol y principales aportes.

% --- TERCER BLOQUE (comentado por defecto) --------------------
%\vspace{-0.1cm}ÁREA DE EXPERIENCIA 3 \hfill (Año inicio - Año fin)\vspace{-0.20cm}
%\noindent\rule{\textwidth}{0.1pt}
%\textbf{(Año) Rol} \textsc{Institución} \hfill País\\
%Descripción breve.

% ==============================================================
%  EDUCACIÓN Y FORMACIÓN ACADÉMICA
% ==============================================================
\vspace{0.2cm}
\textbf{EDUCACIÓN Y FORMACIÓN ACADÉMICA}\vspace{-0.25cm}\\
\noindent\rule{\textwidth}{2pt}

% Ordenar de más reciente a más antiguo.
% Comentar con % los que no sean relevantes para la postulación.
\textbf{(Año - Cursando)} Nombre del programa \hfill \textsc{Institución,} País\\
\textbf{(Año - Año)} Nombre del programa \hfill \textsc{Institución,} País\\
\textbf{(Año - Año)} Nombre del programa \hfill \textsc{Institución,} País\\
% Agrega las líneas que necesites

% ==============================================================
%  RECONOCIMIENTOS (comentar sección completa si no aplica)
% ==============================================================
\vspace{0.2cm}
\textbf{RECONOCIMIENTOS}\vspace{-0.25cm}

\noindent\rule{\textwidth}{1pt}
\textbf{(Año) Nombre del reconocimiento} en \textsc{Institución} \hfill Ciudad, País\\
Descripción breve del reconocimiento y su relevancia.\\
% Agrega más entradas si corresponde

% ==============================================================
%  OTRAS CERTIFICACIONES (comentar sección completa si no aplica)
% ==============================================================
\vspace{0.2cm}
\textbf{OTRAS CERTIFICACIONES}\vspace{-0.25cm}

\noindent\rule{\textwidth}{1pt}
\textbf{(Año)} Nombre del curso o certificación \hfill \textsc{Institución}\\
\textbf{(Año)} Nombre del curso o certificación \hfill \textsc{Institución}\\
% Agrega o comenta según la postulación

\end{document}
